\documentclass[a4paper]{article}

\usepackage[english]{babel}
%\usepackage[utf8x]{inputenc}
\usepackage{amsmath}
\usepackage{graphicx}
\usepackage[colorinlistoftodos]{todonotes}
\usepackage[all]{xy}
\usepackage{dsfont}
\usepackage{amssymb}
\usepackage{bbold}
\usepackage{amsthm}

\usepackage{tikz}
\usepackage{verbatim}
\usetikzlibrary{matrix}
\usetikzlibrary{cd}

\usepackage[autostyle=true]{csquotes}
\usepackage{pdfpages}
\usepackage{hyperref}
\usepackage{fancyhdr}



\newtheorem{theorem}{Theorem}[section]
\newtheorem{proposition}{Proposition}[theorem]
\newtheorem{lemma}{Lemma}[theorem]
\newtheorem{corollary}{Corollary}[theorem]
\newtheorem*{definition}{Definition}


\newcommand{\norm}[1]{\left\lVert#1\right\rVert}

\newcommand{\Rn}{\mathbb{R}^n}


\newcommand{\too}[1]{\overset{#1}{\longrightarrow}}
\newcommand{\hookar}[1]{\overset{#1}{\hookrightarrow}}
\newcommand{\colim}[1]{ \underset{#1}{\texttt{colim}} \text{ } }
\newcommand{\fset}[1]{ \{ #1_0 , \dots , #1_n \} }
\newcommand{\ddt}{\frac{d}{dt}}


\newcommand{\N}{\mathbb{N}}
\newcommand{\Z}{\mathbb{Z}}









\begin{document}









\thispagestyle{empty}
\section*{Preguntas}
\begin{enumerate}


     \item Muestre que la funci\'on 
      $T:\mathbb{R}^3\to P_2[x]$, definida por $$T(a,b,c)=(2a+b)+(a-b)x+ (a-c)x^2$$
     es una transformaci\'on lineal. Encuentre el kernel de $T$, una base y su dimensi\'on ¿Cu\'al es la dimensi\'on de la imagen?
     
    


        
     \item Muestre que la funci\'on 
          $T:\mathbb{R}^3\to P_2[x]$, definida por $$T(a,b,c)=(a-b)+(a-2b)x+ (2b+3c)x^2$$
          es una transformaci\'on lineal. Encuentre el kernel de $T$, una base y su dimensi\'on ¿Cu\'al es la dimensi\'on de la imagen?







\end{enumerate}


\end{document}